\documentclass{article}
\usepackage[T1]{fontenc}
\usepackage{lmodern}
\usepackage{graphicx}
\usepackage{amsmath,amssymb}
\usepackage[a4paper,margin=2cm]{geometry}
\usepackage[absolute,overlay]{textpos}
\setlength{\TPHorizModule}{1mm}
\setlength{\TPVertModule}{1mm}
\usepackage[indonesian]{babel}
\usepackage{hyperref}
\setlength{\emergencystretch}{2em}
% unit: Halaman 56

\begin{document}

% === Halaman 56 (python/native) ===

56

Langkah-langkahnya sbb:

1.

Misalkan panjang kunci yang sudah berhasil dideduksi adalah n. Setiap huruf 

berjarak  n pasti dienkripsi dengan huruf kunci yang sama. Kelompokkan setiap 

huruf berjarak n bersama-sama sehingga kriptanalis memiliki n buah “pesan”, 

masing-masing “pesan” dienkripsi dengan huruf kunci yang sama dengan cipher 

abjad-tunggal (dalam hal ini Caesar cipher). 

2.

Tiap-tiap pesan dari hasil langkah 1 dapat dipecahkan dengan metode analisis 

frekuensi.

3.

Dari hasil langkah 2 kriptanalis dapat menyusun huruf-huruf kunci. Atau, 

kriptanalis dapat menerka kata yang membantu untuk memecahkan cipherteks

\end{document}
